% Options for packages loaded elsewhere
% Options for packages loaded elsewhere
\PassOptionsToPackage{unicode}{hyperref}
\PassOptionsToPackage{hyphens}{url}
\PassOptionsToPackage{dvipsnames,svgnames,x11names}{xcolor}
%
\documentclass[
]{article}
\usepackage{xcolor}
\usepackage{amsmath,amssymb}
\setcounter{secnumdepth}{5}
\usepackage{iftex}
\ifPDFTeX
  \usepackage[T1]{fontenc}
  \usepackage[utf8]{inputenc}
  \usepackage{textcomp} % provide euro and other symbols
\else % if luatex or xetex
  \usepackage{unicode-math} % this also loads fontspec
  \defaultfontfeatures{Scale=MatchLowercase}
  \defaultfontfeatures[\rmfamily]{Ligatures=TeX,Scale=1}
\fi
\usepackage{lmodern}
\ifPDFTeX\else
  % xetex/luatex font selection
\fi
% Use upquote if available, for straight quotes in verbatim environments
\IfFileExists{upquote.sty}{\usepackage{upquote}}{}
\IfFileExists{microtype.sty}{% use microtype if available
  \usepackage[]{microtype}
  \UseMicrotypeSet[protrusion]{basicmath} % disable protrusion for tt fonts
}{}
\makeatletter
\@ifundefined{KOMAClassName}{% if non-KOMA class
  \IfFileExists{parskip.sty}{%
    \usepackage{parskip}
  }{% else
    \setlength{\parindent}{0pt}
    \setlength{\parskip}{6pt plus 2pt minus 1pt}}
}{% if KOMA class
  \KOMAoptions{parskip=half}}
\makeatother
% Make \paragraph and \subparagraph free-standing
\makeatletter
\ifx\paragraph\undefined\else
  \let\oldparagraph\paragraph
  \renewcommand{\paragraph}{
    \@ifstar
      \xxxParagraphStar
      \xxxParagraphNoStar
  }
  \newcommand{\xxxParagraphStar}[1]{\oldparagraph*{#1}\mbox{}}
  \newcommand{\xxxParagraphNoStar}[1]{\oldparagraph{#1}\mbox{}}
\fi
\ifx\subparagraph\undefined\else
  \let\oldsubparagraph\subparagraph
  \renewcommand{\subparagraph}{
    \@ifstar
      \xxxSubParagraphStar
      \xxxSubParagraphNoStar
  }
  \newcommand{\xxxSubParagraphStar}[1]{\oldsubparagraph*{#1}\mbox{}}
  \newcommand{\xxxSubParagraphNoStar}[1]{\oldsubparagraph{#1}\mbox{}}
\fi
\makeatother


\usepackage{longtable,booktabs,array}
\usepackage{calc} % for calculating minipage widths
% Correct order of tables after \paragraph or \subparagraph
\usepackage{etoolbox}
\makeatletter
\patchcmd\longtable{\par}{\if@noskipsec\mbox{}\fi\par}{}{}
\makeatother
% Allow footnotes in longtable head/foot
\IfFileExists{footnotehyper.sty}{\usepackage{footnotehyper}}{\usepackage{footnote}}
\makesavenoteenv{longtable}
\usepackage{graphicx}
\makeatletter
\newsavebox\pandoc@box
\newcommand*\pandocbounded[1]{% scales image to fit in text height/width
  \sbox\pandoc@box{#1}%
  \Gscale@div\@tempa{\textheight}{\dimexpr\ht\pandoc@box+\dp\pandoc@box\relax}%
  \Gscale@div\@tempb{\linewidth}{\wd\pandoc@box}%
  \ifdim\@tempb\p@<\@tempa\p@\let\@tempa\@tempb\fi% select the smaller of both
  \ifdim\@tempa\p@<\p@\scalebox{\@tempa}{\usebox\pandoc@box}%
  \else\usebox{\pandoc@box}%
  \fi%
}
% Set default figure placement to htbp
\def\fps@figure{htbp}
\makeatother


% definitions for citeproc citations
\NewDocumentCommand\citeproctext{}{}
\NewDocumentCommand\citeproc{mm}{%
  \begingroup\def\citeproctext{#2}\cite{#1}\endgroup}
\makeatletter
 % allow citations to break across lines
 \let\@cite@ofmt\@firstofone
 % avoid brackets around text for \cite:
 \def\@biblabel#1{}
 \def\@cite#1#2{{#1\if@tempswa , #2\fi}}
\makeatother
\newlength{\cslhangindent}
\setlength{\cslhangindent}{1.5em}
\newlength{\csllabelwidth}
\setlength{\csllabelwidth}{3em}
\newenvironment{CSLReferences}[2] % #1 hanging-indent, #2 entry-spacing
 {\begin{list}{}{%
  \setlength{\itemindent}{0pt}
  \setlength{\leftmargin}{0pt}
  \setlength{\parsep}{0pt}
  % turn on hanging indent if param 1 is 1
  \ifodd #1
   \setlength{\leftmargin}{\cslhangindent}
   \setlength{\itemindent}{-1\cslhangindent}
  \fi
  % set entry spacing
  \setlength{\itemsep}{#2\baselineskip}}}
 {\end{list}}
\usepackage{calc}
\newcommand{\CSLBlock}[1]{\hfill\break\parbox[t]{\linewidth}{\strut\ignorespaces#1\strut}}
\newcommand{\CSLLeftMargin}[1]{\parbox[t]{\csllabelwidth}{\strut#1\strut}}
\newcommand{\CSLRightInline}[1]{\parbox[t]{\linewidth - \csllabelwidth}{\strut#1\strut}}
\newcommand{\CSLIndent}[1]{\hspace{\cslhangindent}#1}



\setlength{\emergencystretch}{3em} % prevent overfull lines

\providecommand{\tightlist}{%
  \setlength{\itemsep}{0pt}\setlength{\parskip}{0pt}}



 


\makeatletter
\@ifpackageloaded{caption}{}{\usepackage{caption}}
\AtBeginDocument{%
\ifdefined\contentsname
  \renewcommand*\contentsname{Table of contents}
\else
  \newcommand\contentsname{Table of contents}
\fi
\ifdefined\listfigurename
  \renewcommand*\listfigurename{List of Figures}
\else
  \newcommand\listfigurename{List of Figures}
\fi
\ifdefined\listtablename
  \renewcommand*\listtablename{List of Tables}
\else
  \newcommand\listtablename{List of Tables}
\fi
\ifdefined\figurename
  \renewcommand*\figurename{Figure}
\else
  \newcommand\figurename{Figure}
\fi
\ifdefined\tablename
  \renewcommand*\tablename{Table}
\else
  \newcommand\tablename{Table}
\fi
}
\@ifpackageloaded{float}{}{\usepackage{float}}
\floatstyle{ruled}
\@ifundefined{c@chapter}{\newfloat{codelisting}{h}{lop}}{\newfloat{codelisting}{h}{lop}[chapter]}
\floatname{codelisting}{Listing}
\newcommand*\listoflistings{\listof{codelisting}{List of Listings}}
\makeatother
\makeatletter
\makeatother
\makeatletter
\@ifpackageloaded{caption}{}{\usepackage{caption}}
\@ifpackageloaded{subcaption}{}{\usepackage{subcaption}}
\makeatother
\usepackage{bookmark}
\IfFileExists{xurl.sty}{\usepackage{xurl}}{} % add URL line breaks if available
\urlstyle{same}
\hypersetup{
  pdftitle={Information Criteria for Multiply-Imputed Data: A Principled Approach for Complete-Data Replication},
  pdfauthor={Marcus R. Waldman},
  pdfkeywords={multiple imputation, information
criteria, AIC, BIC, model selection, missing data, Q-function},
  colorlinks=true,
  linkcolor={blue},
  filecolor={Maroon},
  citecolor={Blue},
  urlcolor={Blue},
  pdfcreator={LaTeX via pandoc}}


\title{Information Criteria for Multiply-Imputed Data: A Principled
Approach for Complete-Data Replication}
\author{Marcus R. Waldman}
\date{2026-01-11}
\begin{document}
\maketitle
\begin{abstract}
Model selection via information criteria (AIC, BIC) is fundamental to
quantitative research, yet standard approaches break down when data
contain missing values. Applied researchers using multiple imputation
(MI) typically average IC values across imputations---a procedure
lacking theoretical justification---or use methods like AICcd that
target predictive performance rather than what researchers actually
want: replication of the inference they would have made had their data
been complete. This article introduces the \emph{complete-data
replication} principle and derives corrected information criteria that
target this goal. We show that the pooled deviance across imputations
(the Q-function) is biased for the complete-data log-likelihood, with
bias equal to ½tr(RIV), where RIV is the relative increase in variance
from Rubin's combining rules. This leads to corrected formulas---MI-AIC
and MI-BIC---that require only standard MI output and represent minimal
modification to current practice. We clarify why our approach differs
from AICcd (which uses 2tr(RIV) rather than tr(RIV)) by showing they
target different inferential goals. Simulation studies in structural
equation modeling and latent profile analysis demonstrate that the
corrected criteria better replicate complete-data model selection
decisions. Practical guidance for implementation is provided.
\end{abstract}


\section{1. Introduction}\label{sec-intro}

Model selection via information criteria (IC) such as AIC (Akaike,
Hirotugu 1974) and BIC (Schwarz, Gideon 1978) are fundamental to
quantitative research across the behavioral sciences. These criteria are
most frequently used for evaluating nonnested models, providing a
principled basis for balancing fit against complexity. For example, in
mixture modeling, IC guide class enumeration decisions (Nylund,
Asparouhov, Muthén 2007; Masyn 2013), while in structural equation
modeling, they inform choices among alternative factor structures and
tests of measurement invariance (Kuiper 2022; Lin, Huang, Weng 2017;
Vrieze 2012; Huang 2017), as well as decisions about cross-loadings and
residual covariances (Lin, Huang, Weng 2017; Bollen et al. 2014). At the
same time, missing data are nearly ubiquitous in applied research.
Multiple imputation (MI) has emerged as the dominant approach for
handling missing data (Rubin, Donald B. 1987; Schafer and Graham 2002;
Graham 2009), offering a framework that is both principled and
convenient for propagating uncertainty from the imputation process
through to final inferences.

Despite the widespread use of both IC and MI, there is no consensus on
how to properly pool them. Applied researchers and the software they use
typically average IC values across imputations---a procedure lacking
both theoretical justification and a clear inferential target. Although
theoretical work on IC for missing data exists, such as the AIC with a
correction for missing data (AICcd) proposed by Cavanaugh and Shumway
(1998), these methods have not been implemented in standard SEM software
and target predictive performance at the observed-data maximum
likelihood estimate rather than what applied researchers typically seek:
replication of the inference they would have made had their data been
complete. This leaves a fundamental question unresolved: \emph{What
information criteria would researchers have computed had their data been
complete?} Although asymptotic theory establishes that IC will select
the correct model structure as sample size grows (Akaike, Hirotugu 1974;
Schwarz, Gideon 1978; Huang 2017), this provides no guidance for the
finite-sample behavior of IC under MI or when all candidate models are
approximations---the most common situation in practice.

The primary purpose of this article is to introduce the
\emph{complete-data replication} principle and derive information
criteria that answer this question by targeting what AIC and BIC would
have been at the complete-data maximum likelihood estimate. We show that
the pooled deviance across imputations is biased for the complete-data
log-likelihood, with bias equal to ½tr(RIV), where RIV is the relative
increase in variance from Rubin's rules (Rubin, Donald B. 1987).
Although this result builds on the connection between MI and the
Q-function from EM-based approaches, our contribution extends this work
to the MI framework where imputations are drawn from a posterior
distribution rather than a point estimate. This leads to corrected
formulas---MI-AIC and MI-BIC---that require only standard MI output. We
clarify why our approach differs from AICcd by showing they target
different inferential goals, conduct simulation studies in structural
equation modeling and latent profile analysis to demonstrate that the
corrected criteria better replicate complete-data model selection
decisions, and provide practical guidance for implementation.

The next section reviews information criteria and multiple imputation,
establishing notation and the complete-data replication principle.
Section 3 presents the theoretical derivation of the bias in the
Q-function and the corrected IC formulas. Section 4 describes two
simulation studies that evaluate performance when the true model is
among the candidates and when all models are approximations. Section 5
presents an empirical example, and Section 6 concludes with a discussion
of implications and recommendations.

\begin{center}\rule{0.5\linewidth}{0.5pt}\end{center}

\section{2. Background: Model Selection and Missing
Data}\label{sec-background}

\subsection{2.1 AIC \& BIC}\label{aic-bic}

The Akaike Information Criterion (AIC; Akaike, Hirotugu (1974)) and
Bayesian Information Criterion (BIC; Schwarz, Gideon (1978)) are model
selection tools that estimate the expected predictive performance of
candidate models. Both are based on the Kullback-Leibler (KL)
divergence, a measure of the distance between a fitted model and the
true data-generating model (Cavanaugh and Neath 2019). AIC provides an
asymptotically unbiased estimator of the expected KL-divergence, while
BIC approximates the Bayes factor for model comparison under certain
prior specifications.

AIC is defined as:

\[\text{AIC} = -2\ell(\hat{\theta}) + 2q\]

BIC is defined as:

\[\text{BIC} = -2\ell(\hat{\theta}) + q\log(n)\]

where \(\ell(\hat{\theta})\) is the log-likelihood evaluated at the
maximum likelihood estimate (MLE), \(q\) is the number of parameters,
and \(n\) is the sample size. The penalty terms (\(2q\) for AIC,
\(q\log(n)\) for BIC) correct for overfitting by penalizing model
complexity. Lower values indicate better models.

While both criteria target predictive performance, they differ in their
asymptotic properties. AIC optimizes expected prediction error, while
BIC is consistent for selecting the true model when it is among the
candidates (Huang 2017). In latent class and mixture modeling contexts,
BIC tends to outperform AIC in selecting the correct number of classes
(Nylund, Asparouhov, Muthén 2007).

Both AIC and BIC assume complete data when computing the log-likelihood
at the MLE. When data contain missing values, this assumption no longer
holds, motivating the need for modified IC formulas.

\subsection{2.2 Multiple Imputation
Framework}\label{multiple-imputation-framework}

Multiple imputation (MI; Rubin, Donald B. (1987)) is a widely-used
approach to treat missing data that is both principled and convenient.
It is principled because it reflects uncertainty about the missing
values through the imputation distribution and, fundamentally,
marginalizes across the complete-data likelihood to approximate the
observed-data likelihood.\footnote{MI validity requires the missing data
  mechanism to be ignorable (i.e., missing at random, MAR) and the
  parameters of the missingness mechanism to be a priori independent of
  the parameters of the analysis model (Rubin, Donald B. 1987).} It is
convenient because it separates the imputation step from the analysis
step, allowing analysts to use standard complete-data methods after
imputation (Schafer and Graham 2002; Graham 2009). The MI procedure
consists of three steps:

\begin{enumerate}
\def\labelenumi{\arabic{enumi}.}
\tightlist
\item
  For \(m = 1, \ldots, M\): Generate one random imputation of the
  missing values under a specified distribution
\item
  Fit the model of interest to each imputed dataset
\item
  Pool the resulting parameter estimates and inferences
\end{enumerate}

For the theoretical results in this article, we assume imputations are
\emph{Bayesianly proper} (Schafer 1997)---that is, the imputed values
are drawn from a posterior predictive distribution that accounts for
uncertainty in both the missing data and the parameters of the
imputation model. This distinguishes proper MI from simpler imputation
methods that condition on point estimates of parameters.

The Bayesian foundation of multiple imputation is to make inference from
the observed-data posterior distribution,
\(p(\theta | Y_{\text{obs}})\). By generating multiple draws from the
posterior predictive distribution of the missing data and analyzing each
completed dataset, MI approximates this observed-data posterior. The
goal is to obtain parameter estimates and uncertainty quantification
that would arise from directly analyzing the observed-data posterior,
without explicitly deriving its analytical form.

\subsubsection{2.2.1 Rubin's Pooling Rules}\label{rubins-pooling-rules}

Rubin's rules (Rubin, Donald B. 1987) provide the framework for
combining estimates across imputations. These combining rules are
derived under the assumption that the observed-data posterior can be
approximated as a normal distribution, allowing for straightforward
pooling of point estimates and variances. For a scalar parameter with
point estimate \(\hat{\theta}_m\) from imputed dataset \(m\):

\begin{itemize}
\tightlist
\item
  \textbf{Pooled estimate}:
  \(\bar{\theta} = M^{-1}\sum_{m=1}^M \hat{\theta}_m\)
\item
  \textbf{Within-imputation variance}: \(W = M^{-1}\sum_{m=1}^M V_m\),
  where \(V_m\) is the variance estimate from dataset \(m\)
\item
  \textbf{Between-imputation variance}:
  \(B = \frac{1}{M-1}\sum_{m=1}^M(\hat{\theta}_m - \bar{\theta})^2\)
\item
  \textbf{Total variance}: \(T = W + (1 + M^{-1})B\)
\item
  \textbf{Relative Increase in Variance}:
  \(\text{RIV} = (1 + M^{-1})B / W\)
\end{itemize}

The within-imputation variance \(W\) reflects the average uncertainty
within each imputed dataset---the variability that would remain even if
the missing values were known. The between-imputation variance \(B\)
captures the variability across imputations, reflecting uncertainty
about what the missing values actually are. The total variance \(T\)
combines these sources, with the term \((1 + M^{-1})B\) providing a
finite-sample correction for using \(M\) rather than infinite
imputations. The relative increase in variance (RIV) quantifies how much
additional uncertainty is introduced by the missingness, expressed as a
proportion of the within-imputation variance. Conceptually, RIV connects
to the missing information principle: it estimates the ratio of missing
information to observed information,
\(\mathcal{I}_{\text{mis}|\text{obs}} / \mathcal{I}_{\text{obs}}\).

\subsubsection{2.2.2 The Q-Function and MI}\label{the-q-function-and-mi}

The connection between MI and the EM algorithm
{[}(\textbf{dempstera.p?}).;lairdn.m.;rubind.b.MaximumLikelihoodIncomplete1977{]}
provides important theoretical context for our derivation. The EM
algorithm iterates between an E-step (computing the expected
complete-data log-likelihood given current parameter estimates) and an
M-step (maximizing this expectation). The Q-function is the expected
complete-data log-likelihood:

\[Q(\theta|\theta_t) = E[\ell_{\text{com}}(\theta) | Y_{\text{obs}}, \theta_t]\]

where \(\ell_{\text{com}}(\theta)\) is the log-likelihood for the
complete data (observed plus imputed), and the expectation is over the
distribution of the missing data given the observed data and current
parameters \(\theta_t\). In traditional EM, this expectation is
evaluated either analytically or through numerical integration.

Stochastic EM (NEEDS CITATION) provides a bridge between traditional EM
and multiple imputation. Rather than computing the expectation
analytically, Stochastic EM approximates the E-step by drawing samples
from the conditional distribution of the missing data. The Q-function is
then approximated as an average of complete-data log-likelihoods across
these draws. Multiple imputation extends this stochastic approach by
additionally drawing parameters from their posterior distribution,
rather than conditioning on point estimates as in Stochastic EM. This
Bayesian treatment of parameter uncertainty distinguishes
\emph{Bayesianly proper} MI (Schafer 1997) from deterministic or
stochastic EM approaches.

Central to this framework is the \emph{missing information principle}
(Orchard and Woodbury 1972; Schafer 1997), which decomposes the
complete-data observed information matrix (derived from the Hessian) as:

\[\mathcal{I}_{\text{com}} = \mathcal{I}_{\text{obs}} + \mathcal{I}_{\text{mis}|\text{obs}}\]

This decomposition reveals that the observed-data information
(\(\mathcal{I}_{\text{obs}}\)) is less than the complete-data
information (\(\mathcal{I}_{\text{com}}\)) by an amount equal to the
missing information (\(\mathcal{I}_{\text{mis}|\text{obs}}\)). The ratio
\(\mathcal{I}_{\text{mis}|\text{obs}} / \mathcal{I}_{\text{obs}}\)
quantifies the information loss due to missingness, which connects
directly to RIV in Rubin's rules. This article extends theoretical
results developed in the Q-function and EM literature to the Bayesianly
proper MI framework, where the target is replication of complete-data
inference rather than observed-data maximization.

\subsection{2.3 Existing Approaches to IC Under Missing
Data}\label{existing-approaches-to-ic-under-missing-data}

\subsubsection{2.3.1 Ad Hoc Averaging}\label{ad-hoc-averaging}

The most common approach in applied practice is to compute IC for each
imputed dataset separately, then average across imputations:

\[\text{IC}_{\text{adhoc}} = M^{-1}\sum_{m=1}^M \text{IC}_m\]

This approach has intuitive appeal---it treats IC values like parameter
estimates and applies the same averaging used for point estimates. It is
also what standard MI software packages compute by default, making it
the path of least resistance for applied researchers. However, this
procedure lacks theoretical justification. It is unclear what quantity
the averaged IC estimates, and the approach conflates two sources of
variability: the within-imputation uncertainty captured by each
\(\text{IC}_m\) and the between-imputation variability reflecting
missing-data uncertainty. Despite these limitations, ad hoc averaging
remains widespread in practice due to its simplicity and default
availability in software.

\subsubsection{2.3.2 PDIO and AICcd}\label{pdio-and-aiccd}

The first theoretically grounded correction for information criteria
under incomplete data was derived by (\textbf{shimodaira1994?}), who
proposed PDIO (predictive divergence for indirect observation models).
Shimodaira showed that for models where observations are only available
indirectly---including missing data, latent variable models, and mixture
distributions---the AIC penalty should be \(\text{tr}(I_X I_Y^{-1})\)
rather than simply \(q\), where \(I_X\) is the complete-data Fisher
information and \(I_Y\) is the observed-data Fisher information. Using
the missing information principle, this can be rewritten as
\(\text{tr}(I_X I_Y^{-1}) = q + \text{tr}(\text{RIV})\), yielding a
penalty of \(2q + 2\text{tr}(\text{RIV})\) on the AIC scale.

Cavanaugh and Shumway (1998) independently derived an equivalent result,
which they termed AICcd (AIC with correction for incomplete data):

\[\text{AICcd} = -2\bar{Q} + 2q + 2\text{tr}(\text{RIV})\]

where \(\bar{Q}\) is the pooled log-likelihood across imputations. Both
PDIO and AICcd target the expected predictive performance of the model
at the observed-data MLE \(\hat{\theta}_{\text{obs}}\). The penalty term
\(2\text{tr}(\text{RIV})\) accounts for the additional degrees of
freedom consumed by estimating parameters in the presence of missing
data---it is double the trace of the missing information ratio because
it corrects for both the bias from in-sample optimization and the
information loss from missingness. Importantly, PDIO and AICcd answer a
different question than the one posed in this article. While they target
predictive performance given the observed-data structure, our MI-AIC
targets replication of what the researcher would have computed had the
data been complete. Both are theoretically valid for their respective
inferential goals, but they optimize different objectives.

\subsubsection{2.3.3 Other Methods}\label{other-methods}

An alternative strategy is to avoid imputation altogether by using full
information maximum likelihood (FIML), which directly maximizes the
observed-data likelihood and allows standard IC formulas to be applied
(Lai 2021). While FIML and MI are both valid approaches, they involve
different computational strategies and may yield different IC values for
the same model. A key observation across existing methods is that none
explicitly target complete-data replication---the inference a researcher
would have made had there been no missing data. This is the gap our
approach fills.

\begin{center}\rule{0.5\linewidth}{0.5pt}\end{center}

\section{3. Theory: Deriving Corrected IC for Complete-Data
Replication}\label{sec-theory}

\subsection{3.1 Defining the Target}\label{defining-the-target}

When researchers use multiple imputation to handle missing data, a
natural question arises: which model would I have selected had my data
been complete? This framing treats MI as a device to approximate
complete-data inference, allowing researchers to make decisions that
replicate what they would have concluded without missingness. Our target
is therefore \(\text{AIC}_{\text{com}}\), the AIC that would have been
computed had the data been complete:

\[\text{AIC}_{\text{com}} = -2\ell_{\text{com}}(\hat{\theta}_{\text{com}}) + 2q\]

where \(\hat{\theta}_{\text{com}}\) is the MLE under the complete-data
model, \(\ell_{\text{com}}(\cdot)\) is the complete-data log-likelihood,
and \(q\) is the number of parameters.

This target differs fundamentally from alternative approaches. For
example, AICcd (Cavanaugh and Shumway 1998) targets predictive
performance evaluated at the observed-data MLE
\(\hat{\theta}_{\text{obs}}\), reflecting the goal of optimal prediction
given the observed-data structure. In contrast, our goal is to replicate
the inference a researcher would make if the data had no missing
values---answering ``what model would I select with complete data?''
rather than ``what model predicts best given my observed data?'' Both
targets are valid, but they serve different inferential purposes. For
most applications in structural equation modeling and applied
psychological research, complete-data replication is the more natural
choice because researchers typically ask counterfactual questions about
their substantive conclusions, not predictive questions conditional on
the pattern of missingness.

\subsection{3.2 What MI Computes}\label{what-mi-computes}

Multiple imputation provides \(M\) completed datasets
\(Y_1^{\text{com}}, \ldots, Y_M^{\text{com}}\), each consisting of the
observed data \(Y^{\text{obs}}\) combined with imputed values for the
missing data. For each completed dataset, we can compute the
complete-data log-likelihood
\(\ell_{\text{com}}(\theta | Y_m^{\text{com}})\) and fit the model to
obtain parameter estimates. The Q-function, introduced in Section 2.2.2,
is the average of the complete-data log-likelihoods across imputations:

\[\bar{Q}(\theta) = M^{-1}\sum_{m=1}^M \ell_{\text{com}}(\theta | Y_m^{\text{com}})\]

In practice, we evaluate this Q-function at the observed-data MLE
\(\hat{\theta}_{\text{obs}}\)---the parameter estimate obtained by
pooling estimates across imputations using Rubin's rules.

To understand the challenge, consider the parameter landscape. Let
\(\theta^*\) denote the true parameter value. The complete-data MLE
\(\hat{\theta}_{\text{com}}\) is the value that would maximize the
complete-data log-likelihood if we had no missing values:
\(\hat{\theta}_{\text{com}} = \arg\max_\theta \ell_{\text{com}}(\theta | Y^{\text{com}})\).
The observed-data MLE \(\hat{\theta}_{\text{obs}}\) is the value
obtained from the observed data (or equivalently, pooled across
imputations), which maximizes
\(\ell_{\text{obs}}(\theta | Y^{\text{obs}})\). Finally, let
\(\tilde{\phi}\) denote the imputation parameter---in Bayesianly proper
MI, this is a draw from the posterior distribution
\(p(\theta|Y^{\text{obs}})\). This posterior draw distinguishes proper
MI from deterministic approaches: each imputation uses a different
\(\tilde{\phi}^{(m)}\) sampled from the posterior, rather than
conditioning on a single point estimate as in the EM algorithm.

We use \(\tilde{\phi}\) to denote the \textbf{imputation model}
parameter to distinguish it from \(\theta\), the \textbf{analysis model}
parameter. This notational distinction is important because the
imputation model and analysis model need not be the same---a situation
known as \textbf{uncongeniality}
(\textbf{mengMultipleimputationInferencesUncongenial1994?}). While our
derivation assumes congeniality (both models target the same true
parameter \(\theta^*\)), the notation keeps the conceptual distinction
clear: \(\tilde{\phi}\) is drawn from the imputation model's posterior
\(p(\phi|Y^{\text{obs}})\) to generate missing data, while \(\theta\)
parameterizes the model we fit to the completed datasets.

The central question is: what is the bias when we use
\(\bar{Q}(\hat{\theta}_{\text{obs}})\) as a proxy for
\(\ell_{\text{com}}(\hat{\theta}_{\text{com}})\)? We want the
log-likelihood at the complete-data MLE, but MI gives us the Q-function
at the observed-data MLE. The bias decomposes into two terms, which we
derive in the next section.

\subsection{3.3 The Bias Decomposition}\label{the-bias-decomposition}

The bias of \(\bar{Q}\) as an estimate of
\(\ell_{\text{com}}(\hat{\theta}_{\text{com}})\) decomposes as:

\[\text{Bias} = \mathbb{E}[\bar{Q}_{MI}(\hat{\theta}_{\text{obs}})] - \mathbb{E}[\ell_{\text{com}}(\hat{\theta}_{\text{com}})] = \underbrace{\mathbb{E}[Q(\hat{\theta}_{\text{obs}}|\tilde{\phi})] - \mathbb{E}[\ell_{\text{com}}(\hat{\theta}_{\text{obs}})]}_{\text{Term 1: Imputation Bias}} + \underbrace{\mathbb{E}[\ell_{\text{com}}(\hat{\theta}_{\text{obs}})] - \mathbb{E}[\ell_{\text{com}}(\hat{\theta}_{\text{com}})]}_{\text{Term 2: Estimation Mismatch}}\]

The notation \(Q(\hat{\theta}_{\text{obs}}|\tilde{\phi})\) emphasizes
that the Q-function depends on both: (1) the parameter
\(\hat{\theta}_{\text{obs}}\) at which we evaluate the log-likelihood,
and (2) the imputation parameter \(\tilde{\phi}\) used to generate the
imputed values. This distinguishes the MI framework from deterministic
EM, where the Q-function is conditioned on a point estimate rather than
a draw from the posterior.

\subsubsection{3.3.1 Term 1: Imputation
Bias}\label{term-1-imputation-bias}

Term 1 quantifies the bias introduced by imputing missing data with a
parameter \(\tilde{\phi}\) drawn from the posterior distribution
\(p(\theta|Y^{\text{obs}})\) rather than the true parameter
\(\theta^*\). A key insight is that if the imputation parameter were the
true parameter, this term would be exactly zero by the law of iterated
expectations:

\[\mathbb{E}[Q(\hat{\theta}_{\text{obs}}|\theta^*)] = \mathbb{E}[\mathbb{E}_{Y^{\text{mis}}|Y^{\text{obs}},\theta^*}[\ell_{\text{com}}(\hat{\theta}_{\text{obs}})]] = \mathbb{E}[\ell_{\text{com}}(\hat{\theta}_{\text{obs}})]\]

The bias arises precisely because \(\tilde{\phi}\) is sampled from the
posterior and therefore \(\tilde{\phi} \neq \theta^*\).

\textbf{Derivation.} We use the log-likelihood factorization
\(\ell_{\text{com}} = \ell_{\text{obs}} + \ell_{\text{mis}|\text{obs}}\).
Since \(\ell_{\text{obs}}(\hat{\theta}_{\text{obs}})\) does not depend
on \(Y^{\text{mis}}\), the \(\ell_{\text{obs}}\) terms cancel:

\[Q(\hat{\theta}_{\text{obs}}|\tilde{\phi}) - \ell_{\text{com}}(\hat{\theta}_{\text{obs}}) = \mathbb{E}_{Y^{\text{mis}}|Y^{\text{obs}},\tilde{\phi}}[\ell_{\text{mis}|\text{obs}}(\hat{\theta}_{\text{obs}})] - \ell_{\text{mis}|\text{obs}}(\hat{\theta}_{\text{obs}})\]

The bias arises entirely from imputing under \(\tilde{\phi}\) versus the
true conditional distribution under \(\theta^*\). Taking the full
expectation and rearranging nested expectations:

\[\text{Term 1} = \mathbb{E}_{Y^{\text{obs}}}\left[\mathbb{E}_{\tilde{\phi}}[\ell_{\text{mis}|\text{obs}}(\hat{\theta}_{\text{obs}})] - \mathbb{E}_{\theta^*}[\ell_{\text{mis}|\text{obs}}(\hat{\theta}_{\text{obs}})]\right]\]

where
\(\mathbb{E}_{\theta}[\cdot] = \mathbb{E}_{Y^{\text{mis}}|Y^{\text{obs}},\theta}[\cdot]\).
Taylor expanding around \(\tilde{\phi} = \theta^*\) and using properties
of exponential family models, the difference becomes:

\[\mathbb{E}_{\tilde{\phi}}[\ell_{\text{mis}|\text{obs}}(\hat{\theta}_{\text{obs}})] - \mathbb{E}_{\theta^*}[\ell_{\text{mis}|\text{obs}}(\hat{\theta}_{\text{obs}})] \approx (\tilde{\phi} - \theta^*)^T \mathbb{E}_{\theta^*}[\ell_{\text{mis}|\text{obs}}(\hat{\theta}_{\text{obs}}) \cdot S_{\text{mis}|\text{obs}}(\theta^*)]\]

where \(S_{\text{mis}|\text{obs}}(\theta^*)\) is the score function.
Expanding \(\ell_{\text{mis}|\text{obs}}(\hat{\theta}_{\text{obs}})\)
around \(\theta^*\) and applying Gibbs' inequality (the cross-entropy
\(\mathbb{E}_{\theta}[\ell_{\text{mis}|\text{obs}}(\theta^*)]\) is
maximized at \(\theta = \theta^*\)), the term
\(\mathbb{E}_{\theta^*}[\ell_{\text{mis}|\text{obs}}(\theta^*) \cdot S_{\text{mis}|\text{obs}}(\theta^*)]\)
vanishes, leaving:

\[\text{Term 1} \approx \mathbb{E}_{Y^{\text{obs}}}\left[(\tilde{\phi} - \theta^*)^T \mathcal{I}_{\text{mis}|\text{obs}}(\theta^*)(\hat{\theta}_{\text{obs}} - \theta^*)\right]\]

Using the trace identity \(a^T B c = \text{tr}(B c a^T)\) and
recognizing that both estimators are asymptotically unbiased:

\[\text{Term 1} = \text{tr}\left(\mathcal{I}_{\text{mis}|\text{obs}}(\theta^*) \cdot \text{Cov}(\tilde{\phi}, \hat{\theta}_{\text{obs}})\right)\]

\textbf{Simplification under key assumptions.} To simplify
\(\text{Cov}(\tilde{\phi}, \hat{\theta}_{\text{obs}})\), we decompose
the posterior draw as:

\[\tilde{\phi} = \hat{\theta}_{\text{obs}} + \Delta\tilde{\phi}\]

where \(\Delta\tilde{\phi}\) represents the deviation of the posterior
draw from the observed-data MLE. This decomposition requires two
assumptions:

\begin{enumerate}
\def\labelenumi{\arabic{enumi}.}
\item
  \textbf{Congeniality}
  (\textbf{mengMultipleimputationInferencesUncongenial1994?}): The
  imputation model is congenial with the analysis model, meaning both
  target the same parameter: \(\mathbb{E}[\tilde{\phi}] \to \theta^*\)
  as \(n \to \infty\). This ensures that
  \(\mathbb{E}[\Delta\tilde{\phi}] = 0\) asymptotically.
\item
  \textbf{Uncorrelatedness}: The deviation \(\Delta\tilde{\phi}\) is
  uncorrelated with \(\hat{\theta}_{\text{obs}}\):
  \(\text{Cov}(\Delta\tilde{\phi}, \hat{\theta}_{\text{obs}}) = 0\).
  This states that the direction of deviation from the MLE is
  uncorrelated with the MLE itself.
\end{enumerate}

Under these assumptions:

\[\text{Cov}(\tilde{\phi}, \hat{\theta}_{\text{obs}}) = \text{Cov}(\hat{\theta}_{\text{obs}} + \Delta\tilde{\phi}, \hat{\theta}_{\text{obs}}) = \text{Var}(\hat{\theta}_{\text{obs}}) + \text{Cov}(\Delta\tilde{\phi}, \hat{\theta}_{\text{obs}}) = \text{Var}(\hat{\theta}_{\text{obs}}) = \mathcal{I}_{\text{obs}}(\theta^*)^{-1}\]

Therefore:

\[\text{Term 1} = \text{tr}\left(\mathcal{I}_{\text{mis}|\text{obs}}(\theta^*) \cdot \mathcal{I}_{\text{obs}}(\theta^*)^{-1}\right) = \text{tr}(\text{RIV})\]

\textbf{Intuition.} Imputing missing data using parameters drawn from
the posterior creates ``artificial coherence'' between the observed and
imputed portions of the data. Because the posterior concentrates around
values that fit the observed data well, the imputed values tend to align
more closely with the observed data than they would if generated under
the true \(\theta^*\). This spurious coherence inflates the
complete-data log-likelihood.

\textbf{Relationship to prior work.} This result corresponds to the bias
correction derived by (\textbf{shimodaira1994?}) for PDIO and
independently by Cavanaugh and Shumway (1998) for AICcd. However, our
derivation is more general because it applies to the multiple imputation
framework where \(\tilde{\phi}\) is a posterior draw rather than a point
estimate. The \(\text{tr}(\text{RIV})\) result holds under the
congeniality and uncorrelatedness assumptions, regardless of the
specific form of the posterior distribution.

\subsubsection{3.3.2 Term 2: Estimation
Mismatch}\label{term-2-estimation-mismatch}

Term 2 represents the loss in complete-data log-likelihood from
evaluating at \(\hat{\theta}_{\text{obs}}\) (the observed-data MLE)
rather than at \(\hat{\theta}_{\text{com}}\) (the complete-data MLE).
Since \(\hat{\theta}_{\text{com}}\) maximizes
\(\ell_{\text{com}}(\theta)\) by definition, we have
\(\ell_{\text{com}}(\hat{\theta}_{\text{com}}) \geq \ell_{\text{com}}(\hat{\theta}_{\text{obs}})\),
which means Term 2 is non-positive. We now derive its exact value.

\textbf{Taylor expansion.} Expanding
\(\ell_{\text{com}}(\hat{\theta}_{\text{obs}})\) in a second-order
Taylor series around the maximizer \(\hat{\theta}_{\text{com}}\):

\[\ell_{\text{com}}(\hat{\theta}_{\text{obs}}) = \ell_{\text{com}}(\hat{\theta}_{\text{com}}) + \nabla\ell_{\text{com}}(\hat{\theta}_{\text{com}})^T(\hat{\theta}_{\text{obs}} - \hat{\theta}_{\text{com}}) - \frac{1}{2}(\hat{\theta}_{\text{obs}} - \hat{\theta}_{\text{com}})^T \mathcal{J}_{\text{com}}(\hat{\theta}_{\text{com}})(\hat{\theta}_{\text{obs}} - \hat{\theta}_{\text{com}}) + o_p(n^{-3/2})\]

where
\(\mathcal{J}_{\text{com}}(\hat{\theta}_{\text{com}}) = -\nabla^2\ell_{\text{com}}(\theta)|_{\theta=\hat{\theta}_{\text{com}}}\)
is the observed information matrix at \(\hat{\theta}_{\text{com}}\).
Because \(\hat{\theta}_{\text{com}}\) is the MLE, the score vanishes:
\(\nabla\ell_{\text{com}}(\hat{\theta}_{\text{com}}) = 0\). Therefore
the first-order term disappears, leaving:

\[\ell_{\text{com}}(\hat{\theta}_{\text{obs}}) - \ell_{\text{com}}(\hat{\theta}_{\text{com}}) = -\frac{1}{2}(\hat{\theta}_{\text{obs}} - \hat{\theta}_{\text{com}})^T \mathcal{J}_{\text{com}}(\hat{\theta}_{\text{com}})(\hat{\theta}_{\text{obs}} - \hat{\theta}_{\text{com}}) + o_p(n^{-3/2})\]

\textbf{Taking expectations.} Under standard regularity conditions, the
observed information converges in probability to the Fisher information:
\(\mathcal{J}_{\text{com}}(\hat{\theta}_{\text{com}}) \xrightarrow{p} \mathcal{I}_{\text{com}}(\theta^*)\).
Replacing with Fisher information and taking expectations:

\[\text{Term 2} = \mathbb{E}[\ell_{\text{com}}(\hat{\theta}_{\text{obs}})] - \mathbb{E}[\ell_{\text{com}}(\hat{\theta}_{\text{com}})] = -\frac{1}{2}\mathbb{E}\left[(\hat{\theta}_{\text{obs}} - \hat{\theta}_{\text{com}})^T \mathcal{I}_{\text{com}}(\theta^*)(\hat{\theta}_{\text{obs}} - \hat{\theta}_{\text{com}})\right] + O(n^{-3/2})\]

Using the trace identity \(a^T B a = \text{tr}(B a a^T)\) and
recognizing that both MLEs are asymptotically unbiased for \(\theta^*\):

\[\text{Term 2} = -\frac{1}{2}\text{tr}\left(\mathcal{I}_{\text{com}}(\theta^*) \cdot \text{Var}(\hat{\theta}_{\text{obs}} - \hat{\theta}_{\text{com}})\right) + O(n^{-3/2})\]

\textbf{Computing the variance of the difference.} To evaluate this
expression, we need the variance decomposition:

\[\text{Var}(\hat{\theta}_{\text{obs}} - \hat{\theta}_{\text{com}}) = \text{Var}(\hat{\theta}_{\text{obs}}) + \text{Var}(\hat{\theta}_{\text{com}}) - 2\text{Cov}(\hat{\theta}_{\text{obs}}, \hat{\theta}_{\text{com}})\]

The marginal variances are standard:
\(\text{Var}(\hat{\theta}_{\text{obs}}) = \mathcal{I}_{\text{obs}}(\theta^*)^{-1}\)
and
\(\text{Var}(\hat{\theta}_{\text{com}}) = \mathcal{I}_{\text{com}}(\theta^*)^{-1}\).
The key is computing the covariance.

\textbf{Asymptotic linear representation.} By standard MLE theory, to
leading order:

\[\hat{\theta}_{\text{obs}} - \theta^* \approx \mathcal{I}_{\text{obs}}(\theta^*)^{-1} S_{\text{obs}}(\theta^*), \quad \hat{\theta}_{\text{com}} - \theta^* \approx \mathcal{I}_{\text{com}}(\theta^*)^{-1} S_{\text{com}}(\theta^*)\]

where \(S_{\text{obs}}(\theta^*)\) and \(S_{\text{com}}(\theta^*)\) are
the score functions evaluated at the true parameter. Therefore:

\[\text{Cov}(\hat{\theta}_{\text{obs}}, \hat{\theta}_{\text{com}}) = \mathcal{I}_{\text{obs}}^{-1} \text{Cov}(S_{\text{obs}}(\theta^*), S_{\text{com}}(\theta^*)) \mathcal{I}_{\text{com}}^{-1}\]

\textbf{Score covariance.} Since
\(S_{\text{com}}(\theta^*) = S_{\text{obs}}(\theta^*) + S_{\text{mis}|\text{obs}}(\theta^*)\):

\[\text{Cov}(S_{\text{obs}}(\theta^*), S_{\text{com}}(\theta^*)) = \text{Var}(S_{\text{obs}}(\theta^*)) + \text{Cov}(S_{\text{obs}}(\theta^*), S_{\text{mis}|\text{obs}}(\theta^*))\]

The key observation is that
\(\text{Cov}(S_{\text{obs}}, S_{\text{mis}|\text{obs}}) = 0\). To see
this, condition on \(Y^{\text{obs}}\): \(S_{\text{obs}}(\theta^*)\) is a
function of \(Y^{\text{obs}}\) only (so it is constant given
\(Y^{\text{obs}}\)), while
\(\mathbb{E}[S_{\text{mis}|\text{obs}}(\theta^*) | Y^{\text{obs}}] = 0\)
by the score property. By the law of total covariance:

\[\text{Cov}(S_{\text{obs}}, S_{\text{mis}|\text{obs}}) = \mathbb{E}[\text{Cov}(S_{\text{obs}}, S_{\text{mis}|\text{obs}}|Y^{\text{obs}})] + \text{Cov}(\mathbb{E}[S_{\text{obs}}|Y^{\text{obs}}], \mathbb{E}[S_{\text{mis}|\text{obs}}|Y^{\text{obs}}]) = 0 + 0 = 0\]

Therefore:

\[\text{Cov}(S_{\text{obs}}(\theta^*), S_{\text{com}}(\theta^*)) = \text{Var}(S_{\text{obs}}(\theta^*)) = \mathcal{I}_{\text{obs}}(\theta^*)\]

\textbf{MLE covariance.} Substituting back:

\[\text{Cov}(\hat{\theta}_{\text{obs}}, \hat{\theta}_{\text{com}}) = \mathcal{I}_{\text{obs}}^{-1} \mathcal{I}_{\text{obs}} \mathcal{I}_{\text{com}}^{-1} = \mathcal{I}_{\text{com}}^{-1} = \text{Var}(\hat{\theta}_{\text{com}})\]

This reveals that the covariance between the two MLEs equals the
variance of the complete-data MLE. Intuitively, the observed-data MLE
``fully overlaps'' with the complete-data MLE in the information it
provides.

\textbf{Variance of the difference.} Using the variance decomposition:

\[\text{Var}(\hat{\theta}_{\text{obs}} - \hat{\theta}_{\text{com}}) = \mathcal{I}_{\text{obs}}^{-1} + \mathcal{I}_{\text{com}}^{-1} - 2\mathcal{I}_{\text{com}}^{-1} = \mathcal{I}_{\text{obs}}^{-1} - \mathcal{I}_{\text{com}}^{-1}\]

This is positive semi-definite, as expected, since
\(\mathcal{I}_{\text{com}}^{-1} \preceq \mathcal{I}_{\text{obs}}^{-1}\)
(the complete-data MLE has lower variance than the observed-data MLE).

\textbf{Substituting back.} Returning to Term 2:

\[\text{Term 2} = -\frac{1}{2}\text{tr}\left(\mathcal{I}_{\text{com}} \cdot \left(\mathcal{I}_{\text{obs}}^{-1} - \mathcal{I}_{\text{com}}^{-1}\right)\right) = -\frac{1}{2}\text{tr}\left(\mathcal{I}_{\text{com}} \mathcal{I}_{\text{obs}}^{-1}\right) + \frac{1}{2}\text{tr}\left(\mathcal{I}_{\text{com}} \mathcal{I}_{\text{com}}^{-1}\right)\]

The second term simplifies to
\(\frac{1}{2}\text{tr}(I_q) = \frac{q}{2}\), where \(q\) is the number
of parameters. For the first term, use the missing information principle
\(\mathcal{I}_{\text{com}} = \mathcal{I}_{\text{obs}} + \mathcal{I}_{\text{mis}|\text{obs}}\):

\[\text{tr}\left(\mathcal{I}_{\text{com}} \mathcal{I}_{\text{obs}}^{-1}\right) = \text{tr}\left((\mathcal{I}_{\text{obs}} + \mathcal{I}_{\text{mis}|\text{obs}}) \mathcal{I}_{\text{obs}}^{-1}\right) = \text{tr}(I_q) + \text{tr}\left(\mathcal{I}_{\text{mis}|\text{obs}} \mathcal{I}_{\text{obs}}^{-1}\right) = q + \text{tr}(\text{RIV})\]

Therefore:

\[\text{Term 2} = -\frac{1}{2}\left(q + \text{tr}(\text{RIV})\right) + \frac{q}{2} = -\frac{q}{2} - \frac{1}{2}\text{tr}(\text{RIV}) + \frac{q}{2} = -\frac{1}{2}\text{tr}(\text{RIV})\]

\textbf{Interpretation.} Evaluating the complete-data log-likelihood at
the observed-data MLE \(\hat{\theta}_{\text{obs}}\) rather than at the
complete-data MLE \(\hat{\theta}_{\text{com}}\) incurs an expected loss
of \(\frac{1}{2}\text{tr}(\text{RIV})\). This loss partially offsets the
positive bias from imputation (Term 1), which was
\(\text{tr}(\text{RIV})\). The factor of \(\frac{1}{2}\) arises from the
quadratic curvature of the log-likelihood surface.

\subsubsection{3.3.3 Total Bias}\label{total-bias}

Combining the two terms:

\[\text{Bias} = \text{Term 1} + \text{Term 2} = \text{tr}(\text{RIV}) - \frac{1}{2}\text{tr}(\text{RIV}) = \frac{1}{2}\text{tr}(\text{RIV})\]

The total bias is positive but only half of the imputation bias. This
partial cancellation has a clear interpretation. Imputation with
estimated parameters creates artificial coherence, inflating the
Q-function by \(\text{tr}(\text{RIV})\) (Term 1). However, because we
evaluate at the observed-data MLE \(\hat{\theta}_{\text{obs}}\) rather
than the complete-data MLE \(\hat{\theta}_{\text{com}}\), we incur a
loss of \(\frac{1}{2}\text{tr}(\text{RIV})\) (Term 2) due to not being
at the maximizer. The net effect is that half of the imputation bias
remains.

The factor of \(\frac{1}{2}\) in the net bias arises from the quadratic
curvature of the log-likelihood surface. The imputation bias increases
linearly with the deviation \((\tilde{\phi} - \theta^*)\), but the
evaluation mismatch involves a quadratic form in
\((\hat{\theta}_{\text{obs}} - \hat{\theta}_{\text{com}})\), yielding
the \(\frac{1}{2}\) factor in the second-order Taylor expansion.

The total bias is directly proportional to the relative increase in
variance (RIV), which quantifies the fraction of information lost due to
missingness. When there is little missing data, RIV is small and the
bias correction is negligible. When missingness is substantial, RIV is
large and the bias correction becomes important for accurate model
selection.

\subsection{3.4 Corrected IC Formulas}\label{corrected-ic-formulas}

The bias of \(\frac{1}{2}\text{tr}(\text{RIV})\) in the log-likelihood
translates directly to a correction in the AIC penalty. Recall that AIC
is defined as \(-2\ell(\hat{\theta}) + 2q\), where the factor of 2
converts the log-likelihood to the deviance scale. Since the Q-function
\(\bar{Q}\) has bias \(\frac{1}{2}\text{tr}(\text{RIV})\) for the
complete-data log-likelihood, the deviance \(-2\bar{Q}\) has bias
\(-2 \times \frac{1}{2}\text{tr}(\text{RIV}) = -\text{tr}(\text{RIV})\).
To correct this bias, we add \(\text{tr}(\text{RIV})\) to the penalty:

\[\text{MI-AIC} = -2\bar{Q} + 2q + \text{tr}(\text{RIV})\]

This formula targets \(\text{AIC}_{\text{com}}\), the AIC that would
have been computed with complete data.

The same correction applies to BIC and other information criteria. BIC
has penalty structure \(q\log(n)\) instead of \(2q\), but the bias still
affects the deviance term \(-2\ell\), not the penalty structure.
Therefore:

\[\text{MI-BIC} = -2\bar{Q} + q\log(n) + \text{tr}(\text{RIV})\]

More generally, for any IC of the form \(-2\ell + f(q, n)\), the
MI-corrected version is \(-2\bar{Q} + f(q,n) + \text{tr}(\text{RIV})\).

\textbf{Practical computation.} All quantities in these formulas are
directly computable from standard MI output. The Q-function \(\bar{Q}\)
is the average log-likelihood across imputations, and the correction
term uses Rubin's within-imputation (\(W\)) and between-imputation
(\(B\)) covariance matrices:

\[\text{tr}(\text{RIV}) = \text{tr}\left((1 + M^{-1})W^{-1}B\right)\]

This computation requires only the covariance matrices that MI software
routinely provides, making the correction straightforward to implement
in practice.

\textbf{Regularity conditions.} These formulas assume standard
asymptotic regularity: the model is correctly specified, the sample size
\(n\) is large enough for Fisher information to be well-defined, and the
number of imputations \(M\) is sufficient to make Monte Carlo error
negligible. The within-imputation covariance matrix \(W\) must be
invertible; in high-dimensional settings where \(W\) is near-singular,
regularization (e.g., ridge-type adjustments) may be necessary before
computing \(W^{-1}B\).

\subsection{3.5 Comparison to PDIO and
AICcd}\label{comparison-to-pdio-and-aiccd}

PDIO (\textbf{shimodaira1994?}) and AICcd (Cavanaugh and Shumway 1998)
represent earlier theoretically grounded approaches to IC under missing
data. Both arrive at the same penalty:

\[\text{PDIO} = \text{AICcd} = -2\bar{Q} + 2q + 2\text{tr}(\text{RIV})\]

The penalty is \(2\text{tr}(\text{RIV})\)---twice our
correction---because these methods target a different inferential goal.
While we target \(\text{AIC}_{\text{com}}\) (complete-data replication),
PDIO and AICcd target predictive performance evaluated at the
observed-data MLE \(\hat{\theta}_{\text{obs}}\).

The difference can be understood through our bias decomposition.
Shimodaira's derivation shows that the total penalty for indirect
observation models should be
\(\text{tr}(I_X I_Y^{-1}) = q + \text{tr}(\text{RIV})\), which on the
AIC scale (multiplied by 2) gives \(2q + 2\text{tr}(\text{RIV})\). This
corresponds to our Term 1 plus the base penalty \(q\). Our contribution
is to further decompose the bias when targeting
\(\hat{\theta}_{\text{com}}\) rather than \(\hat{\theta}_{\text{obs}}\):
Term 2 (estimation mismatch) contributes
\(-\frac{1}{2}\text{tr}(\text{RIV})\), partially offsetting Term 1 and
yielding a net correction of only \(\text{tr}(\text{RIV})\). Both
derivations are correct for their respective targets.

Table 1 summarizes the comparison across methods:

\begin{longtable}[]{@{}
  >{\raggedright\arraybackslash}p{(\linewidth - 6\tabcolsep) * \real{0.2500}}
  >{\raggedright\arraybackslash}p{(\linewidth - 6\tabcolsep) * \real{0.2500}}
  >{\raggedright\arraybackslash}p{(\linewidth - 6\tabcolsep) * \real{0.2500}}
  >{\raggedright\arraybackslash}p{(\linewidth - 6\tabcolsep) * \real{0.2500}}@{}}
\toprule\noalign{}
\begin{minipage}[b]{\linewidth}\raggedright
Method
\end{minipage} & \begin{minipage}[b]{\linewidth}\raggedright
Penalty
\end{minipage} & \begin{minipage}[b]{\linewidth}\raggedright
Target Estimand
\end{minipage} & \begin{minipage}[b]{\linewidth}\raggedright
Use When\ldots{}
\end{minipage} \\
\midrule\noalign{}
\endhead
\bottomrule\noalign{}
\endlastfoot
Complete-data AIC & \(2q\) & Prediction (complete data) & No missing
data \\
MI-AIC (ours) & \(2q + \text{tr}(\text{RIV})\) & Complete-data
replication & Asking ``what model would I select with complete
data?'' \\
PDIO/AICcd & \(2q + 2\text{tr}(\text{RIV})\) & Prediction at
\(\hat{\theta}_{\text{obs}}\) & Asking ``what model predicts best given
observed structure?'' \\
Ad hoc average & \(2q\) (incorrect) & Unclear & Should not be
used---lacks theory \\
\end{longtable}

\textbf{Which method to use?} The choice depends on the research
question. If the goal is to replicate the model selection decision one
would make with complete data---the most common scenario in applied
structural equation modeling and psychological research---use MI-AIC or
MI-BIC. If the goal is to optimize predictive performance conditional on
the observed-data structure, use PDIO or AICcd. The ad hoc approach of
simply averaging IC across imputations lacks theoretical justification
and should be avoided, as it conflates within- and between-imputation
variability without a clear target estimand.

All three corrected methods (MI-AIC, PDIO, and AICcd) are asymptotically
valid and computationally straightforward, requiring only standard MI
output. The key is choosing the method that matches the inferential
question of interest.

\begin{center}\rule{0.5\linewidth}{0.5pt}\end{center}

\section{4. Simulation Studies}\label{sec-simulation}

\subsection{4.1 Overview}\label{overview}

Three simulation studies validate the theory and demonstrate practical
performance:

\begin{enumerate}
\def\labelenumi{\arabic{enumi}.}
\tightlist
\item
  \textbf{Study 1}: A simplified pedagogical simulation to build
  intuition about the bias formula and its key assumptions
\item
  \textbf{Study 2}: Replication of Bollen et al.'s (2014) SEM design
  with missing data and MI
\item
  \textbf{Study 3}: Latent profile analysis (LPA) simulation for mixture
  modeling context
\end{enumerate}

\subsection{4.2 Study 1: Illustrating the Bias Formula and Key
Assumptions}\label{study-1-illustrating-the-bias-formula-and-key-assumptions}

\subsubsection{4.2.1 Purpose}\label{purpose}

Before evaluating model selection performance in realistic settings, we
use a simplified simulation to build intuition about the bias formula
derived in Section 3. This study addresses three questions:

\begin{enumerate}
\def\labelenumi{\arabic{enumi}.}
\tightlist
\item
  Does the bias formula (Bias = ½tr(RIV)) hold across different missing
  data mechanisms?
\item
  What role does congeniality play in the bias formula?
\item
  How do inclusive imputation models affect tr(RIV) and the penalty
  term?
\end{enumerate}

The key insight we aim to demonstrate is that \emph{congeniality}---not
ignorability---is the critical assumption for the bias formula. The
missing data mechanism (MCAR, MAR, MNAR) affects whether point estimates
are unbiased for the true parameter \(\theta^*\), but it does not affect
whether the relationship between the Q-function and complete-data
log-likelihood follows our derivation. In contrast, when congeniality is
violated, the bias formula breaks down regardless of the missing
mechanism.

\subsubsection{4.2.2 Design}\label{design}

\textbf{Data generating mechanism.} We simulate a simple three-variable
mediation system with variables X, M, and Y. The true data generating
mechanism is full mediation:

\begin{itemize}
\tightlist
\item
  \(X \sim N(0, 1)\)
\item
  \(M = \beta_{XM} X + \epsilon_M\), where
  \(\epsilon_M \sim N(0, \sigma^2_M)\)
\item
  \(Y = \beta_{MY} M + \epsilon_Y\), where
  \(\epsilon_Y \sim N(0, \sigma^2_Y)\)
\end{itemize}

Note that X has no direct effect on Y in the true model. We set
\(\beta_{XM} = 0.5\), \(\beta_{MY} = 0.5\), and
\(\sigma^2_M = \sigma^2_Y = 0.75\) to yield standardized path
coefficients of approximately 0.5.

\textbf{Sample size and missingness.} Each simulated dataset has
\(N = 100\) observations. We induce substantial missingness:
approximately 50\% of observations are missing M, Y, or both. This high
rate of missingness ensures that the bias and penalty terms are large
enough to observe clearly.

\textbf{Conditions.} We cross two factors in a 3 × 3 factorial design:

\emph{Missing data mechanism (3 levels):}

\begin{itemize}
\tightlist
\item
  \textbf{MCAR}: Missingness is independent of all variables. Each
  observation has a 50\% probability of being selected for missingness,
  with M, Y, or both missing at random among those selected.
\item
  \textbf{MAR}: Missingness depends on X (the fully observed variable).
  Observations with higher X values have higher probability of
  missingness on M and Y.
\item
  \textbf{MNAR}: Missingness depends on M or Y themselves. Observations
  with higher M values have higher probability of M being missing;
  similarly for Y.
\end{itemize}

\emph{Imputation model (3 levels):}

\begin{itemize}
\tightlist
\item
  \textbf{Matched}: The imputation model matches the true data
  generating mechanism (X → M → Y). This is congenial with both analysis
  models.
\item
  \textbf{Saturated}: The imputation model assumes a fully saturated
  covariance structure among X, M, and Y. This is congenial but more
  inclusive than the true model, resulting in wider posterior variance.
\item
  \textbf{Uncongenial}: The imputation model assumes X → M and X → Y
  with M and Y conditionally independent given X. This is uncongenial
  because it contradicts the true M → Y relationship.
\end{itemize}

\textbf{Analysis models.} For each imputed dataset, we fit two
structural equation models:

\begin{itemize}
\tightlist
\item
  \textbf{Model A (Full mediation)}: X → M → Y with no direct X → Y
  path. This is the true model.
\item
  \textbf{Model B (Partial mediation)}: X → M → Y with an additional
  direct X → Y path. This model is overparameterized relative to the
  true DGM.
\end{itemize}

\textbf{Number of imputations and replications.} We generate \(M = 20\)
imputations per dataset and conduct 1,000 replications per condition.

\subsubsection{4.2.3 Outcomes}\label{outcomes}

For each replication, we compute:

\begin{enumerate}
\def\labelenumi{\arabic{enumi}.}
\item
  \textbf{Empirical bias}:
  \(\bar{Q} - \ell_{\text{com}}(\hat{\theta}_{\text{com}})\), comparing
  the MI-pooled Q-function to the complete-data log-likelihood computed
  from the full data (before missingness was induced).
\item
  \textbf{Theoretical bias}: \(\frac{1}{2}\text{tr}(\text{RIV})\)
  computed from Rubin's combining rules.
\item
  \textbf{Agreement}: The correspondence between empirical and
  theoretical bias across conditions.
\end{enumerate}

We examine these outcomes separately for each analysis model (full
mediation vs.~partial mediation) to understand how model complexity
interacts with the imputation model.

\subsubsection{4.2.4 Expected Results}\label{expected-results}

\textbf{Result 1: The missing mechanism does not affect the bias
formula.}

Across MCAR, MAR, and MNAR conditions---holding the imputation model
fixed at either Matched or Saturated (both congenial)---the empirical
bias should approximately equal ½tr(RIV). The missing mechanism affects
whether point estimates \(\hat{\theta}_{\text{obs}}\) are unbiased for
\(\theta^*\), but not whether the bias decomposition from Section 3
holds. This demonstrates that ignorability is not the relevant
assumption for our bias formula.

\textbf{Result 2: Congeniality is essential for the bias formula.}

Under the uncongenial imputation model (X → M, X → Y), the empirical
bias will \emph{not} equal ½tr(RIV). Because this imputation model
assumes M and Y are conditionally independent given X---contradicting
the true M → Y relationship---the congeniality assumption is violated.
Specifically, \(\mathbb{E}[\tilde{\phi}] \not\to \theta^*\), so the
decomposition
\(\tilde{\phi} = \hat{\theta}_{\text{obs}} + \Delta\tilde{\phi}\) with
\(\mathbb{E}[\Delta\tilde{\phi}] = 0\) fails. This result will hold
regardless of the missing mechanism (MCAR, MAR, or MNAR), reinforcing
that congeniality---not ignorability---is the critical assumption.

\textbf{Result 3: Inclusive imputation inflates tr(RIV) but preserves
the bias formula.}

The saturated imputation model, while congenial, produces systematically
larger tr(RIV) than the matched imputation model. This occurs because
the saturated model has wider posterior variance---it estimates more
parameters and thus has greater uncertainty. In terms of our
decomposition, \(\text{Var}(\Delta\tilde{\phi}) > 0\) for the saturated
model, inflating \(\text{Var}(\tilde{\phi})\).

Critically, even with this inflated variance:

\begin{itemize}
\tightlist
\item
  The bias formula still holds: empirical bias ≈ ½tr(RIV)
\item
  The larger tr(RIV) yields a correspondingly larger penalty term
\end{itemize}

This confirms the theoretical result from Section 3.3.1: inclusive
imputation models do not distort the bias formula because
\(\text{Var}(\Delta\tilde{\phi})\) does not enter
\(\text{Cov}(\tilde{\phi}, \hat{\theta}_{\text{obs}})\).

\textbf{Result 4: The penalty term self-corrects for imputation model
complexity.}

The inflated tr(RIV) under saturated imputation has an important
implication: without proper penalization, the larger Q-function values
from inclusive imputation could spuriously favor more complex analysis
models. However, because the penalty term scales with tr(RIV), this
tendency is offset. The MI-AIC and MI-BIC formulas automatically adjust
for the additional uncertainty introduced by inclusive imputation.

This foreshadows the role of proper penalization in Studies 2 and 3,
where we evaluate actual model selection performance.

\subsubsection{4.2.5 Summary}\label{summary}

Study 1 establishes the empirical validity of the theoretical framework
from Section 3:

\begin{enumerate}
\def\labelenumi{\arabic{enumi}.}
\tightlist
\item
  The bias = ½tr(RIV) relationship is robust to the missing data
  mechanism
\item
  Congeniality is the critical assumption---violations cause the formula
  to fail
\item
  Inclusive (but congenial) imputation models inflate tr(RIV), but the
  penalty self-corrects
\item
  The distinction between ignorability (for valid point estimates) and
  congeniality (for the bias formula) is empirically confirmed
\end{enumerate}

These results provide the foundation for interpreting Studies 2 and 3,
which evaluate model selection performance in more realistic settings.

\subsection{4.3 Study 2: Structural Equation
Modeling}\label{study-2-structural-equation-modeling}

{[}Replication of Bollen et al.~(2014) SEM design with missing data and
MI --- to be added{]}

\subsection{4.4 Study 3: Latent Profile
Analysis}\label{study-3-latent-profile-analysis}

{[}LPA simulation for mixture modeling context --- to be added{]}

\begin{center}\rule{0.5\linewidth}{0.5pt}\end{center}

\section{5. Empirical Example}\label{sec-empirical}

{[}Empirical example to be added{]}

\begin{center}\rule{0.5\linewidth}{0.5pt}\end{center}

\section{6. Discussion}\label{sec-discussion}

{[}Discussion to be added{]}

\begin{center}\rule{0.5\linewidth}{0.5pt}\end{center}

\section{Appendix A: Detailed Derivation of Term 1 (Imputation
Bias)}\label{appendix-a}

This appendix provides a complete derivation of the imputation bias term
presented in Section 3.3.1. We show that:

\[\text{Term 1} = \mathbb{E}[Q(\hat{\theta}_{\text{obs}}|\tilde{\phi})] - \mathbb{E}[\ell_{\text{com}}(\hat{\theta}_{\text{obs}})] = \text{tr}\left(\mathcal{I}_{\text{mis}|\text{obs}}(\theta^*) \cdot \text{Cov}(\tilde{\phi}, \hat{\theta}_{\text{obs}})\right)\]

and under the congeniality and uncorrelatedness assumptions, this
simplifies to \(\text{tr}(\text{RIV})\).

\subsection{A.1 Proof Strategy}\label{a.1-proof-strategy}

We use the log-likelihood factorization
\(\ell_{\text{com}} = \ell_{\text{obs}} + \ell_{\text{mis}|\text{obs}}\).
The Q-function can be written as:

\[Q(\hat{\theta}_{\text{obs}}|\tilde{\phi}) = \ell_{\text{obs}}(\hat{\theta}_{\text{obs}}) + \mathbb{E}_{Y^{\text{mis}}|Y^{\text{obs}},\tilde{\phi}}[\ell_{\text{mis}|\text{obs}}(\hat{\theta}_{\text{obs}})]\]

while the complete-data log-likelihood is:

\[\ell_{\text{com}}(\hat{\theta}_{\text{obs}}) = \ell_{\text{obs}}(\hat{\theta}_{\text{obs}}) + \ell_{\text{mis}|\text{obs}}(\hat{\theta}_{\text{obs}})\]

The \(\ell_{\text{obs}}(\hat{\theta}_{\text{obs}})\) terms are identical
and cancel, leaving:

\[Q(\hat{\theta}_{\text{obs}}|\tilde{\phi}) - \ell_{\text{com}}(\hat{\theta}_{\text{obs}}) = \mathbb{E}_{Y^{\text{mis}}|Y^{\text{obs}},\tilde{\phi}}[\ell_{\text{mis}|\text{obs}}(\hat{\theta}_{\text{obs}})] - \ell_{\text{mis}|\text{obs}}(\hat{\theta}_{\text{obs}})\]

Therefore:

\[\text{Term 1} = \mathbb{E}\left[\mathbb{E}_{Y^{\text{mis}}|Y^{\text{obs}},\tilde{\phi}}[\ell_{\text{mis}|\text{obs}}(\hat{\theta}_{\text{obs}})] - \ell_{\text{mis}|\text{obs}}(\hat{\theta}_{\text{obs}})\right]\]

where the outer expectation is over
\((Y^{\text{obs}}, Y^{\text{mis}}) \sim p(Y|\theta^*)\).

\subsection{A.2 Simplifying the Nested
Expectations}\label{a.2-simplifying-the-nested-expectations}

Write the full expectation explicitly:

\[\text{Term 1} = \mathbb{E}_{Y^{\text{obs}}|\theta^*}\left[\mathbb{E}_{Y^{\text{mis}}|Y^{\text{obs}},\theta^*}\left[\mathbb{E}_{Y^{\text{mis}}|Y^{\text{obs}},\tilde{\phi}}[\ell_{\text{mis}|\text{obs}}(\hat{\theta}_{\text{obs}})] - \ell_{\text{mis}|\text{obs}}(\hat{\theta}_{\text{obs}})\right]\right]\]

Note that
\(\hat{\theta}_{\text{obs}} = \hat{\theta}_{\text{obs}}(Y^{\text{obs}})\)
and \(\tilde{\phi} = \tilde{\phi}(Y^{\text{obs}})\) are functions of
\(Y^{\text{obs}}\) only.

Since
\(\mathbb{E}_{Y^{\text{mis}}|Y^{\text{obs}},\tilde{\phi}}[\ell_{\text{mis}|\text{obs}}(\hat{\theta}_{\text{obs}})]\)
is a function of \(Y^{\text{obs}}\) only (we've integrated out
\(Y^{\text{mis}}\)), the middle expectation
\(\mathbb{E}_{Y^{\text{mis}}|Y^{\text{obs}},\theta^*}[\cdot]\) acts only
on the second term:

\[\text{Term 1} = \mathbb{E}_{Y^{\text{obs}}|\theta^*}\left[\mathbb{E}_{\tilde{\phi}}[\ell_{\text{mis}|\text{obs}}(\hat{\theta}_{\text{obs}})] - \mathbb{E}_{\theta^*}[\ell_{\text{mis}|\text{obs}}(\hat{\theta}_{\text{obs}})]\right]\]

where we use the shorthand
\(\mathbb{E}_{\theta}[\cdot] = \mathbb{E}_{Y^{\text{mis}}|Y^{\text{obs}},\theta}[\cdot]\).

This is the difference in expected log-likelihood when imputing under
\(\tilde{\phi}\) versus under the true \(\theta^*\).

\subsection{A.3 Taylor Expansion of the
Expectation}\label{a.3-taylor-expansion-of-the-expectation}

We expand \(\mathbb{E}_{\tilde{\phi}}[g]\) around
\(\tilde{\phi} = \theta^*\) for any function \(g(Y^{\text{mis}})\):

\[\mathbb{E}_{\tilde{\phi}}[g] - \mathbb{E}_{\theta^*}[g] \approx (\tilde{\phi} - \theta^*)^T \left.\frac{\partial}{\partial \theta}\mathbb{E}_{\theta}[g]\right|_{\theta=\theta^*}\]

For exponential family models, this derivative can be expressed using
the score function. The score of the conditional distribution is:

\[S_{\text{mis}|\text{obs}}(\theta) = \frac{\partial}{\partial \theta} \ell_{\text{mis}|\text{obs}}(\theta) = \frac{\partial}{\partial \theta} \log p(Y^{\text{mis}}|Y^{\text{obs}}, \theta)\]

A fundamental property of score functions is that the expected score at
the true parameter equals zero:

\[\mathbb{E}_{\theta^*}[S_{\text{mis}|\text{obs}}(\theta^*)] = 0\]

This follows from the fact that
\(\int p(Y^{\text{mis}}|Y^{\text{obs}}, \theta^*) \, dY^{\text{mis}} = 1\),
and differentiating both sides with respect to \(\theta\) gives zero.

For exponential family models, the derivative of the expectation equals:

\[\left.\frac{\partial}{\partial \theta}\mathbb{E}_{\theta}[g]\right|_{\theta=\theta^*} = \mathbb{E}_{\theta^*}[g \cdot S_{\text{mis}|\text{obs}}(\theta^*)]\]

This is the covariance between \(g\) and the score function (since
\(\mathbb{E}[S] = 0\)).

\subsection{A.4 Applying to Our
Setting}\label{a.4-applying-to-our-setting}

With \(g = \ell_{\text{mis}|\text{obs}}(\hat{\theta}_{\text{obs}})\):

\[\mathbb{E}_{\tilde{\phi}}[\ell_{\text{mis}|\text{obs}}(\hat{\theta}_{\text{obs}})] - \mathbb{E}_{\theta^*}[\ell_{\text{mis}|\text{obs}}(\hat{\theta}_{\text{obs}})] \approx (\tilde{\phi} - \theta^*)^T \mathbb{E}_{\theta^*}[\ell_{\text{mis}|\text{obs}}(\hat{\theta}_{\text{obs}}) \cdot S_{\text{mis}|\text{obs}}(\theta^*)]\]

\subsection{A.5 Expanding the
Log-Likelihood}\label{a.5-expanding-the-log-likelihood}

To evaluate
\(\mathbb{E}_{\theta^*}[\ell_{\text{mis}|\text{obs}}(\hat{\theta}_{\text{obs}}) \cdot S_{\text{mis}|\text{obs}}(\theta^*)]\),
we Taylor expand
\(\ell_{\text{mis}|\text{obs}}(\hat{\theta}_{\text{obs}})\) around
\(\theta^*\):

\[\ell_{\text{mis}|\text{obs}}(\hat{\theta}_{\text{obs}}) \approx \ell_{\text{mis}|\text{obs}}(\theta^*) + (\hat{\theta}_{\text{obs}} - \theta^*)^T S_{\text{mis}|\text{obs}}(\theta^*)\]

Substituting:

\[\mathbb{E}_{\theta^*}[\ell_{\text{mis}|\text{obs}}(\hat{\theta}_{\text{obs}}) \cdot S_{\text{mis}|\text{obs}}(\theta^*)] \approx \mathbb{E}_{\theta^*}[\ell_{\text{mis}|\text{obs}}(\theta^*) \cdot S_{\text{mis}|\text{obs}}(\theta^*)] + (\hat{\theta}_{\text{obs}} - \theta^*)^T \mathbb{E}_{\theta^*}[S_{\text{mis}|\text{obs}}(\theta^*) S_{\text{mis}|\text{obs}}(\theta^*)^T]\]

The second term involves the Fisher information matrix. A fundamental
result in likelihood theory is that the variance of the score equals the
Fisher information:

\[\mathbb{E}_{\theta^*}[S_{\text{mis}|\text{obs}}(\theta^*) S_{\text{mis}|\text{obs}}(\theta^*)^T] = \mathcal{I}_{\text{mis}|\text{obs}}(\theta^*)\]

Therefore:

\[\mathbb{E}_{\theta^*}[\ell_{\text{mis}|\text{obs}}(\hat{\theta}_{\text{obs}}) \cdot S_{\text{mis}|\text{obs}}(\theta^*)] \approx \mathbb{E}_{\theta^*}[\ell_{\text{mis}|\text{obs}}(\theta^*) \cdot S_{\text{mis}|\text{obs}}(\theta^*)] + \mathcal{I}_{\text{mis}|\text{obs}}(\theta^*)(\hat{\theta}_{\text{obs}} - \theta^*)\]

\subsection{A.6 The Key Cancellation (Gibbs'
Inequality)}\label{a.6-the-key-cancellation-gibbs-inequality}

We now show that
\(\mathbb{E}_{\theta^*}[\ell_{\text{mis}|\text{obs}}(\theta^*) \cdot S_{\text{mis}|\text{obs}}(\theta^*)] = 0\).

This term can be recognized as the derivative of a cross-entropy:

\[\mathbb{E}_{\theta^*}[\ell_{\text{mis}|\text{obs}}(\theta^*) \cdot S_{\text{mis}|\text{obs}}(\theta^*)] = \left.\frac{\partial}{\partial \theta}\mathbb{E}_{\theta}[\ell_{\text{mis}|\text{obs}}(\theta^*)]\right|_{\theta=\theta^*}\]

The function
\(\mathbb{E}_{\theta}[\ell_{\text{mis}|\text{obs}}(\theta^*)]\) is a
cross-entropy:

\[\mathbb{E}_{\theta}[\ell_{\text{mis}|\text{obs}}(\theta^*)] = \int p(Y^{\text{mis}}|Y^{\text{obs}},\theta) \log p(Y^{\text{mis}}|Y^{\text{obs}},\theta^*) \, dY^{\text{mis}}\]

By \textbf{Gibbs' inequality} (equivalently, the non-negativity of
KL-divergence), this cross-entropy is maximized when
\(\theta = \theta^*\). This is because:

\[D_{KL}(p_{\theta^*} \| p_{\theta^*}) = 0 \leq D_{KL}(p_\theta \| p_{\theta^*}) \quad \text{for all } \theta\]

where \(D_{KL}(p \| q) = \mathbb{E}_p[\log p - \log q] \geq 0\) with
equality iff \(p = q\).

At the maximum \(\theta = \theta^*\), the derivative with respect to
\(\theta\) must vanish:

\[\mathbb{E}_{\theta^*}[\ell_{\text{mis}|\text{obs}}(\theta^*) \cdot S_{\text{mis}|\text{obs}}(\theta^*)] = 0\]

\subsection{A.7 Combining Results}\label{a.7-combining-results}

From Sections A.4--A.6:

\[\mathbb{E}_{\tilde{\phi}}[\ell_{\text{mis}|\text{obs}}(\hat{\theta}_{\text{obs}})] - \mathbb{E}_{\theta^*}[\ell_{\text{mis}|\text{obs}}(\hat{\theta}_{\text{obs}})] \approx (\tilde{\phi} - \theta^*)^T \mathcal{I}_{\text{mis}|\text{obs}}(\theta^*)(\hat{\theta}_{\text{obs}} - \theta^*)\]

Taking the expectation over \(Y^{\text{obs}}\):

\[\text{Term 1} = \mathbb{E}_{Y^{\text{obs}}}\left[(\tilde{\phi} - \theta^*)^T \mathcal{I}_{\text{mis}|\text{obs}}(\theta^*)(\hat{\theta}_{\text{obs}} - \theta^*)\right]\]

\subsection{A.8 Converting to Trace
Form}\label{a.8-converting-to-trace-form}

Using the trace identity \(a^T B c = \text{tr}(B c a^T)\):

\[\text{Term 1} = \text{tr}\left(\mathcal{I}_{\text{mis}|\text{obs}}(\theta^*) \cdot \mathbb{E}[(\hat{\theta}_{\text{obs}} - \theta^*)(\tilde{\phi} - \theta^*)^T]\right)\]

Since both \(\hat{\theta}_{\text{obs}}\) and \(\tilde{\phi}\) are
asymptotically unbiased for \(\theta^*\), the expectation equals the
covariance:

\[\text{Term 1} = \text{tr}\left(\mathcal{I}_{\text{mis}|\text{obs}}(\theta^*) \cdot \text{Cov}(\tilde{\phi}, \hat{\theta}_{\text{obs}})\right)\]

This is our main result in general form.

\subsection{A.9 Simplification Under Key
Assumptions}\label{a.9-simplification-under-key-assumptions}

To simplify \(\text{Cov}(\tilde{\phi}, \hat{\theta}_{\text{obs}})\), we
decompose the posterior draw as:

\[\tilde{\phi} = \hat{\theta}_{\text{obs}} + \Delta\tilde{\phi}\]

where \(\Delta\tilde{\phi}\) represents the deviation of the posterior
draw from the observed-data MLE.

\textbf{Assumption 1 (Congeniality):} The imputation model is congenial
with the analysis model
(\textbf{mengMultipleimputationInferencesUncongenial1994?}), meaning
both target the same parameter:
\(\mathbb{E}[\tilde{\phi}] \to \theta^*\) as \(n \to \infty\). This
ensures that \(\mathbb{E}[\Delta\tilde{\phi}] = 0\) asymptotically.

\textbf{Assumption 2 (Uncorrelatedness):} The deviation
\(\Delta\tilde{\phi}\) is uncorrelated with
\(\hat{\theta}_{\text{obs}}\):
\(\text{Cov}(\Delta\tilde{\phi}, \hat{\theta}_{\text{obs}}) = 0\).

Under these assumptions:

\[\text{Cov}(\tilde{\phi}, \hat{\theta}_{\text{obs}}) = \text{Cov}(\hat{\theta}_{\text{obs}} + \Delta\tilde{\phi}, \hat{\theta}_{\text{obs}}) = \text{Var}(\hat{\theta}_{\text{obs}}) + \text{Cov}(\Delta\tilde{\phi}, \hat{\theta}_{\text{obs}}) = \text{Var}(\hat{\theta}_{\text{obs}}) = \mathcal{I}_{\text{obs}}(\theta^*)^{-1}\]

Therefore:

\[\text{Term 1} = \text{tr}\left(\mathcal{I}_{\text{mis}|\text{obs}}(\theta^*) \cdot \mathcal{I}_{\text{obs}}(\theta^*)^{-1}\right) = \text{tr}(\text{RIV})\]

where the last equality uses the definition
\(\text{RIV} = \mathcal{I}_{\text{mis}|\text{obs}} \cdot \mathcal{I}_{\text{obs}}^{-1}\).

\subsection{A.10 Summary}\label{a.10-summary}

The imputation bias equals \(\text{tr}(\text{RIV})\) under two
assumptions: (1) the imputation model is congenial with the analysis
model, and (2) the deviation of the posterior draw from the MLE is
uncorrelated with the MLE. These assumptions are satisfied by standard
Bayesianly proper multiple imputation procedures. Importantly, the
result holds regardless of the variance of the posterior draws---wider
posterior variance (as in ``fully inclusive'' imputation models) does
not affect the bias correction.

\begin{center}\rule{0.5\linewidth}{0.5pt}\end{center}

\section{References}\label{references}

{[}References will be compiled from references.bib{]}

\phantomsection\label{refs}
\begin{CSLReferences}{1}{0}
\bibitem[\citeproctext]{ref-akaikehirotuguNewLookStatistical1974}
Akaike, Hirotugu. 1974. {``A {New Look} at the {Statistical Model
Identification}.''} \url{https://doi.org/10.1109/TAC.1974.1100705}.

\bibitem[\citeproctext]{ref-bollenBICAlternativeBayesian2014}
Bollen, Kenneth A., Jeffrey J. Harden, Surajit Ray, and Jane Zavisca.
2014. {``{BIC} and {Alternative Bayesian Information Criteria} in the
{Selection} of {Structural Equation Models}.''} \emph{Structural
Equation Modeling: A Multidisciplinary Journal} 21 (1): 1--19.
\url{https://doi.org/10.1080/10705511.2014.856691}.

\bibitem[\citeproctext]{ref-cavanaughAkaikeInformationCriterion2019}
Cavanaugh, Joseph E., and Andrew A. Neath. 2019. {``The {Akaike}
Information Criterion: {Background}, Derivation, Properties,
Application, Interpretation, and Refinements.''} \emph{WIREs
Computational Statistics} 11 (3): e1460.
\url{https://doi.org/10.1002/wics.1460}.

\bibitem[\citeproctext]{ref-cavanaughAkaikeInformationCriterion1998}
Cavanaugh, Joseph E., and Robert H. Shumway. 1998. {``An {Akaike}
Information Criterion for Model Selection in the Presence of Incomplete
Data.''} \emph{Journal of Statistical Planning and Inference} 67:
45--65. \url{https://doi.org/10.1016/s0378-3758(97)00115-8}.

\bibitem[\citeproctext]{ref-grahamMissingDataAnalysis2009}
Graham. 2009. {``Missing {Data Analysis}: {Making It Work} in the {Real
World}.''} \emph{Annual Review of Psychology}.
\url{https://doi.org/10.1146/annurev.psych.58.110405.085530}.

\bibitem[\citeproctext]{ref-huangAsymptoticsAICBIC2017}
Huang. 2017. {``Asymptotics of {AIC}, {BIC}, and {RMSEA} for {Model
Selection} in {Structural Equation Modeling}.''} \emph{Psychometrika}.
\url{https://doi.org/10.1007/s11336-017-9572-y}.

\bibitem[\citeproctext]{ref-kuiperAICtypeTheoryBasedModel2022}
Kuiper. 2022. {``{AIC-type Theory-Based Model Selection} for {Structural
Equation Models}.''} \emph{Structural Equation Modeling: A
Multidisciplinary Journal}.
\url{https://doi.org/10.1080/10705511.2020.1836967}.

\bibitem[\citeproctext]{ref-laiUsingInformationCriteria2021}
Lai, Keke. 2021. {``Using {Information Criteria Under Missing Data}:
{Full Information Maximum Likelihood Versus Two-Stage Estimation}.''}
\emph{Structural Equation Modeling: A Multidisciplinary Journal} 28 (2):
278--91. \url{https://doi.org/10.1080/10705511.2020.1780925}.

\bibitem[\citeproctext]{ref-linhuangwengSelectingPathModels2017}
Lin, Huang, Weng. 2017. {``Selecting {Path Models} in {SEM}: {A
Comparison} of {Model Selection Criteria}.''} \emph{Structural Equation
Modeling: A Multidisciplinary Journal}.
\url{https://doi.org/10.1080/10705511.2017.1363652}.

\bibitem[\citeproctext]{ref-masynLatentClassAnalysis2013a}
Masyn. 2013. {``Latent {Class Analysis} and {Finite Mixture
Modeling}.''} \emph{Oxford Handbooks Online}.
\url{https://doi.org/10.1093/oxfordhb/9780199934898.013.0025}.

\bibitem[\citeproctext]{ref-nylundasparouhovmuthenDecidingNumberClasses2007}
Nylund, Asparouhov, Muthén. 2007. {``Deciding on the {Number} of
{Classes} in {Latent Class Analysis} and {Growth Mixture Modeling}: {A
Monte Carlo Simulation Study}.''} \emph{Structural Equation Modeling: A
Multidisciplinary Journal}.
\url{https://doi.org/10.1080/10705510701575396}.

\bibitem[\citeproctext]{ref-orchardMissingInformationPrinciple1972a}
Orchard, Terence, and Max A. Woodbury. 1972. {``A Missing Information
Principle: Theory and Applications.''} \emph{Proc. Sixth Berkeley Symp.
On Math. Statist. And} Prob. (Vol.1): 697--715.

\bibitem[\citeproctext]{ref-rubindonaldb.MultipleImputationNonresponse1987}
Rubin, Donald B. 1987. {``Multiple {Imputation} for {Nonresponse} in
{Surveys}.''}

\bibitem[\citeproctext]{ref-schaferAnalysisIncompleteMultivariate1997a}
Schafer, Joseph L. 1997. \emph{Analysis of Incomplete Multivariate
Data}. CRC press.

\bibitem[\citeproctext]{ref-schaferMissingDataOur2002}
Schafer, Joseph L, and John W Graham. 2002. {``Missing Data: Our View of
the State of the Art.''} \emph{Psychological Methods} 7 (2): 147.
\url{https://doi.org/10.1037/1082-989x.7.2.147}.

\bibitem[\citeproctext]{ref-schwarzgideonEstimatingDimensionModel1978}
Schwarz, Gideon. 1978. {``Estimating the {Dimension} of a {Model}.''}
\url{https://doi.org/10.1214/aos/1176344136}.

\bibitem[\citeproctext]{ref-vriezeModelSelectionPsychological2012}
Vrieze. 2012. {``Model Selection and Psychological Theory: {A}
Discussion of the Differences Between the {Akaike} Information Criterion
({AIC}) and the {Bayesian} Information Criterion ({BIC}).''}
\emph{Psychological Methods}. \url{https://doi.org/10.1037/a0027127}.

\end{CSLReferences}




\end{document}
